\section{Main Theorem}
\label{sec:main-theorem}

\begin{theorem}[Conditional deterministic dimension amplification]
\label{thm:main}
Under Assumptions~\ref{assump:source-regime},
\ref{assump:universal-sgd-success}, and
\ref{assump:tie-resolved-confident-map},
\[
\operatorname{dc}(\mathcal H)\leq 7TSd.
\]
The conclusion is deterministic and gives exact tie-resolved representation on
every point of $\mathcal X$.  It is nonasymptotic and uses the fixed finite
horizon $T$.  The exposed variables in the rate are $S,T,d$; there is no
hidden constant, and the numerical factor $7$ is independent of
$n,\mathcal H,L,(n_i)_{i=0}^L,S,\eta,T,\varepsilon,d$, and $\mathcal P$.
The probability and expectation in the assumptions are used only to prove
deterministic structural and existence statements; no confidence parameter,
approximation error, margin, or norm change occurs in the conclusion.
\end{theorem}

The theorem is conditional on the target-independent confident-map premise.
It does not assert that exact SGD success by itself produces that premise or a
polynomial value of $d$.

\begin{corollary}[Polynomial specialization under a separate dimension bound]
\label{cor:main-polynomial-specialization}
Under the assumptions of Theorem~\ref{thm:main}, suppose a separate result
establishes an explicit polynomial expression $p(S,T)$ satisfying
$d\leq p(S,T)$ at the current parameters, with all coefficients and auxiliary
dependence exposed and with no hidden dependence on $n$ or $\eta$.  Then
\[
\operatorname{dc}(\mathcal H)\leq7TSp(S,T).
\]
This is a deterministic, fixed-horizon, exact tie-resolved specialization and
does not claim that the additional inequality follows from SGD.
\end{corollary}

\begin{proof}
Theorem~\ref{thm:main} gives
$\operatorname{dc}(\mathcal H)\leq7TSd$.  Assumption
\ref{assump:source-regime} and the positive widths give $T,S\geq1$, while
$d\leq p(S,T)$ gives $p(S,T)-d\geq0$.  Therefore
\[
7TSp(S,T)-7TSd=7TS\bigl(p(S,T)-d\bigr)\geq0,
\]
and the displayed conclusion follows.  This scalar substitution changes no
probability mode, horizon, representation map, score, or tie convention.  If
$d=0$, the exact dimension-zero representation proved below remains the
baseline; if $p(S,T)=d$, the bridge is equality.
\end{proof}
